% Insertion-ready section.  Requires amsmath, amssymb, amsthm, booktabs, array.
\section{The fifteen $p=1201$ selector rays inside
$\operatorname{Herm}_3(\mathbb H)$}
\label{sec:fifteen-selector-jordan}

The common integer $15$ in the complete $p=1201$ selector census and in
$\dim_{\mathbb R}\operatorname{Herm}_3(\mathbb H)=15$ admits an exact test.
The test gives a nontrivial morphism, but it rules out the interpretation of
the selector states as fifteen independent Jordan-algebra directions.

\subsection{Marked selectors and their prime families}

Fix $p_0=1201$.  A marked state in the complete census is a tuple
\[
 s=(\chi,R,h,r,t),\qquad \chi\in\{E,M\},\qquad
 p_0=4rth-R,\qquad \gcd(r,t)=1,
\]
with positive integral $R,h,r,t$.  Its channel gate is
\begin{equation}
 \begin{array}{ll}
 E:&R\mid4r^2h+1,\\
 M:&R\mid r+t.
 \end{array}
 \label{eq:fifteen-gates}
\end{equation}
In the exterior channel the first gate is equivalent to $R\mid p_0r+t$,
because
\[
 p_0r+t=t(4r^2h+1)-Rr.
\]
Put
\begin{equation}
(u,a,v)=(hr^2,hrt,ht^2),\qquad uv=a^2,\qquad R=4a-p_0.
 \label{eq:fifteen-uav}
\end{equation}
If $\kappa_E=(p_0r+t)/R$ and $m=(r+t)/R$, the ordered denominator
maps are
\begin{align}
 D_E(s)&=(hrt,\ ht\kappa_E,\ p_0hr\kappa_E),\label{eq:fifteen-DE}\\
 D_M(s)&=(hrt,\ p_0hrm,\ p_0htm).\label{eq:fifteen-DM}
\end{align}
Clearing the full common denominator proves
\[
 \sum_{d\in D_\chi(s)}\frac1d=\frac4{p_0}.
\]

The fifteen states, with orientation and scale retained, are as follows.
The integer $L$ is the least prime increment that also preserves
$p\equiv361\pmod{840}$.
\begin{center}
\scriptsize
\begin{tabular}{@{}lccrcrrrrc@{}}
\toprule
tag&$\chi$&$R$&$h$&$(r,t)$&$u$&$a$&$v$&$L$&Koide $+$\\
\midrule
E23a&E&23&17&$(1,18)$&17&306&5508&57960&yes\\
M23a&M&23&3&$(6,17)$&108&306&867&14280&no\\
M23b&M&23&3&$(17,6)$&867&306&108&14280&no\\
E23b&E&23&34&$(9,1)$&2754&306&34&57960&no\\
M31a&M&31&2&$(1,154)$&2&308&47432&9240&yes\\
E31&E&31&77&$(4,1)$&1232&308&77&52080&no\\
M31b&M&31&2&$(154,1)$&47432&308&2&9240&no\\
M39a&M&39&2&$(1,155)$&2&310&48050&26040&yes\\
E39&E&39&62&$(5,1)$&1550&310&62&10920&no\\
M39b&M&39&2&$(155,1)$&48050&310&2&26040&no\\
M47a&M&47&1&$(8,39)$&64&312&1521&43680&no\\
M47b&M&47&1&$(39,8)$&1521&312&64&43680&no\\
E71&E&71&53&$(1,6)$&53&318&1908&59640&no\\
E339&E&339&11&$(35,1)$&13475&385&11&94920&no\\
E1043&E&1043&17&$(33,1)$&18513&561&17&1376760&no\\
\bottomrule
\end{tabular}
\end{center}

For an exterior row define
\[
 h_k=h_0+Rk,\qquad p_k=p_0+4rtRk;
\]
for a middle row define
\[
 h_k=h_0+k,\qquad p_k=p_0+4rtk.
\]
The corresponding gate in \eqref{eq:fifteen-gates} is preserved.  The
exterior prime class modulo $4rtR$ and the middle prime class modulo $4rt$
are reduced.  Their least intersections with $361\pmod{840}$ are
\begin{equation}
 L_s=
 \begin{cases}
  \operatorname{lcm}(4rtR,840),&\chi=E,\\
  \operatorname{lcm}(4rt,840),&\chi=M.
 \end{cases}
 \label{eq:fifteen-hard-L}
\end{equation}
Thus
\begin{equation}
 p_n=1201+L_sn,\qquad h_n=h_0+\frac{L_s}{4rt}n
 \label{eq:fifteen-hard-family}
\end{equation}
retains the same marked ray, and Dirichlet's theorem gives infinitely many
prime $p_n$ in every row.  The fifteen marked rows give eleven unordered
denominator triples at $p_0$ and ten distinct residue classes
$1201\pmod{L_s}$.  These are different quotients: middle reversal produces
the first collapse, while E23a and E23b have distinct rays but the same
$L_s=57960$.

\subsection{The exact Koide cone map}

Define
\begin{equation}
 \Gamma(s)=\begin{pmatrix}u&a\\a&v\end{pmatrix}
 =h\binom rt(r\ t).
 \label{eq:fifteen-Gamma}
\end{equation}
This matrix is positive semidefinite of rank one and has Smith normal form
$\operatorname{diag}(h,0)$.  The Lorentz-coordinate map
\begin{equation}
 \mathcal L(\Gamma)=
 \left(\frac{u+v}{2},a,\frac{u-v}{2}\right)=(H,X,Y)
 \label{eq:fifteen-Lorentz}
\end{equation}
satisfies $H^2-X^2-Y^2=uv-a^2=0$.

Let
\[
 F(q_1,q_2,q_3)=
 \left(\frac{q_1+q_2+q_3}{\sqrt3},
       \frac{q_1-q_3}{\sqrt2},
       \frac{q_1+q_3-2q_2}{\sqrt6}\right).
\]
This is orthogonal.  Its inverse at \eqref{eq:fifteen-Lorentz} is
\begin{align}
 q_1&=\frac{u+v}{2\sqrt3}+\frac a{\sqrt2}
                         +\frac{u-v}{2\sqrt6},\nonumber\\
 q_2&=\frac{u+v}{2\sqrt3}-\frac{u-v}{\sqrt6},\label{eq:fifteen-Koide}\\
 q_3&=\frac{u+v}{2\sqrt3}-\frac a{\sqrt2}
                         +\frac{u-v}{2\sqrt6}.\nonumber
\end{align}
For
\[
 Q_K(q)=\frac23(q_1+q_2+q_3)^2-(q_1^2+q_2^2+q_3^2)
\]
orthogonality gives $Q_K(q)=H^2-X^2-Y^2=0$.  Hence
\eqref{eq:fifteen-Koide} is an exact coefficient-preserving map over
$\mathbb Q(\sqrt2,\sqrt3)$, not a numerical fit.  In the displayed
orientation exactly E23a, M31a, and M39a lie in the positive Koide chamber.

\subsection{Embedding and rank obstruction in the quaternionic Jordan algebra}

Let $J=\operatorname{Herm}_3(\mathbb H)$ with
$A\circ B=(AB+BA)/2$, and put $e=\operatorname{diag}(1,1,0)$.  Define
\begin{equation}
 \iota:\operatorname{Sym}_2(\mathbb R)\longrightarrow J,
 \qquad
 \iota\begin{pmatrix}u&a\\a&v\end{pmatrix}
 =\begin{pmatrix}u&a&0\\a&v&0\\0&0&0\end{pmatrix}.
 \label{eq:fifteen-iota}
\end{equation}
Block multiplication proves
$\iota(A\circ B)=\iota(A)\circ\iota(B)$.  Thus $\iota$ is an injective
real-linear Jordan homomorphism with image
\[
 C_{\mathbb R}:=\iota(\operatorname{Sym}_2(\mathbb R))
 =eJe\cap\operatorname{Herm}_3(\mathbb R),
 \qquad \dim_{\mathbb R}C_{\mathbb R}=3.
\]
The full Peirce corner is
$eJe\cong\operatorname{Herm}_2(\mathbb H)$ of real dimension $6$;
$C_{\mathbb R}$ is its proper real-symmetric Jordan subalgebra.
It preserves trace, trace pairing, positivity, and ordinary rank.  In the
Moore convention
\[
 \begin{pmatrix}d_1&q_3&\bar q_2\\
 \bar q_3&d_2&q_1\\q_2&\bar q_1&d_3\end{pmatrix},
\]
the only nonzero coordinates of $\iota\Gamma(s)$ are
$d_1=u,d_2=v,q_3=a\in\mathbb R$.  Its Moore cubic is zero.

Projectivizing the scale gives the idempotent
\begin{equation}
 P_{r,t}=\frac1{r^2+t^2}
 \begin{pmatrix}r^2&rt&0\\rt&t^2&0\\0&0&0\end{pmatrix},
 \qquad P_{r,t}^2=P_{r,t},\quad\operatorname{tr}P_{r,t}=1.
 \label{eq:fifteen-projector}
\end{equation}
It represents the quaternionic line $[r:t:0]_{\mathbb H}$.  The fifteen
rows have pairwise distinct ratios $[r:t]$, so the composition
\begin{equation}
 \Sigma_{1201}\hookrightarrow\mathbb{RP}^1\hookrightarrow\mathbb{HP}^2,
 \qquad s\longmapsto[r:t]\longmapsto[r:t:0]_{\mathbb H}
 \label{eq:fifteen-RP1-HP2}
\end{equation}
is injective.

\begin{theorem}[Exact span and Jordan closure]
\label{thm:fifteen-span}
The fifteen matrices $\iota\Gamma(s)$ are fifteen distinct projective
rank-one points, but
\[
 \operatorname{span}_{\mathbb R}\{\iota\Gamma(s):s\in\Sigma_{1201}\}
 =\langle\iota\Gamma(\Sigma_{1201})\rangle_{\rm Jordan}=C_{\mathbb R}.
\]
In particular, both spaces have real dimension $3$, not $15$.
\end{theorem}
\begin{proof}
In coordinates
$\bigl(\begin{smallmatrix}x&y\\y&z\end{smallmatrix}\bigr)
\leftrightarrow(x,y,z)$, the projective image of $(r,t)$ is
$(r^2,rt,t^2)$.  The rows E31, E39, and E339 give
\[
 \begin{pmatrix}16&4&1\\25&5&1\\1225&35&1\end{pmatrix},
 \qquad \det=-930\ne0.
\]
Hence the images span the whole three-dimensional real-symmetric corner
$C_{\mathbb R}$.  The space $C_{\mathbb R}$ is a Jordan subalgebra by the
block-multiplication identity above, so the generated Jordan algebra is
contained in $C_{\mathbb R}$.  A Jordan subalgebra is a real vector space;
because it contains the fifteen generators, it contains their span
$C_{\mathbb R}$.  Both containments are equalities.
\end{proof}

The number fifteen is not invariant even after the hard residue class is held
fixed.  Complete divisor enumeration gives the following exact comparison:
\[
\begin{array}{c|rrrr}
 p&1201&4561&8761&9601\\ \hline
 \#E&7&26&16&17\\
 \#M&8&28&20&12\\
 \#(E\sqcup M)&15&54&36&29\\
 \#\{\text{unordered denominator triples}\}&11&40&26&23
\end{array}
\]
All four displayed primes satisfy $p\equiv361\pmod{840}$.  Thus neither the
hard residue class nor the selector construction singles out a
fifteen-element family; the equality
$\#\Sigma_{1201}=\dim_{\mathbb R}\operatorname{Herm}_3(\mathbb H)=15$ is a
feature of the individual prime $1201$, not an invariant of the arithmetic
family.  The deterministic enumeration and primality checks are included in
the accompanying verification script.

At fixed $p_0$ the unprojectivized embedding preserves the markings.  From
$\iota\Gamma$ one recovers $(u,a,v)$ and then
\[
 h=\gcd(u,a,v),\quad r=\sqrt{u/h},\quad t=\sqrt{v/h},\quad
 R=4a-1201,\quad j=(R-3)/4.
\]
The two channel gates cannot overlap in this census: overlap would imply
$R\mid1200$, and none of $23,31,39,47,71,339,1043$ divides $1200$.
Projectivization in \eqref{eq:fifteen-projector} loses $h$, and without the
fixed prime the residual formula $R=4a-p$ is unavailable.

\subsection{Symmetry and incidence obstructions}

Normalized orientation reversal fixes the seven exterior states and swaps
the four middle pairs, so its cycle type is $1^7 2^4$.  The transvection
\[
 \tau(x_1,x_2,x_3,x_4)=(x_1+x_2,x_2,x_3,x_4)
\]
on $\mathbb F_2^4\setminus\{0\}$ has the same cycle type.  For example, with
integers $1,\ldots,15$ encoding the four binary coordinates from the low bit,
the following map is $C_2$-equivariant:
\[
\begin{array}{c|rrrrrrr}
s&E23a&E23b&E31&E39&E71&E339&E1043\\\hline
\phi(s)&1&4&5&8&9&12&13
\end{array}
\]
\[
\begin{array}{c|rrrr}
s&M23a&M23b&M31a&M31b\\\hline
\phi(s)&2&3&6&7
\end{array}
\qquad
\begin{array}{c|rrrr}
s&M39a&M39b&M47a&M47b\\\hline
\phi(s)&10&11&14&15
\end{array}
\]
This proves only an equivalence of $C_2$-sets.  It is noncanonical: a set
with this cycle type has $7!2^4 4!=1{,}935{,}360$ automorphisms commuting
with the involution.

It cannot be upgraded to the natural projective incidence.  The selector
points lie on the nonsingular conic $xz-y^2=0$ in
$\mathbb P^2(\mathbb R)$, so no three distinct selector images are collinear.
By contrast, $PG(3,2)$ has $35$ three-point lines.  Nor does normalized
reversal extend linearly on the canonical real-symmetric corner: it fixes
E31, E39, and E339, which span $C_{\mathbb R}$, so such a linear map would be
the identity; but it must exchange four nonidentical middle pairs.

The other elementary count comparisons do not supply a bridge.  Type $A_3$
has $12$ oriented roots, $6$ unoriented root lines, and $4+6+4=14$ nonzero
minuscule weights.  The rank-one idempotents of $J$ form the real
eight-dimensional continuum $\mathbb{HP}^2$, while the present images occupy
the proper copy $\mathbb{RP}^1$ in \eqref{eq:fifteen-RP1-HP2}.

\subsection{Conclusion and next exact test}

The current structural statement is therefore
\[
 \boxed{
 \begin{gathered}
 \text{fifteen marked prime-carrying selector rays}\ 
 \hookrightarrow\mathbb{RP}^1\hookrightarrow\mathbb{HP}^2,\\
 \dim\operatorname{span}=\dim\operatorname{Jordan\ closure}=3.
 \end{gathered}}
\]
The equality between the number of marked rays and $\dim_{\mathbb R}J$ is
not explained by the defined morphism.

A decisive higher-rung test would construct, from the arithmetic residual and
channel markings, an explicit element $g_s$ of a named quaternionic structure
group and study
\[
 g_s\iota(\Gamma(s))g_s^*.
\]
One must then compute the span and Jordan closure over complete multi-prime
censuses and prove equivariance with the selector involutions.  If all images
remain in one common three-dimensional real-symmetric Jordan corner, the
obstruction persists; if the markings canonically move them through enough
conjugate corners and the closure is all of $J$, that would be a genuine
fifteen-dimensional bridge.  No such
marking-to-conjugation map is presently defined.
